\scanpage{91}
\mspage{45}
$\mathcal E_U$ sur $\mathcal E_S$ \unsure{ou|où} $\unsure{=}$ de $\mathcal E_{\widetilde U}$ sur $\unsure{\mathcal E_{\widetilde\sigma}|\mathcal E_{\widetilde S}}$, définit néanmoins, \deleted{\unsure{le}} le sous-groupe $\mathcal E_{\widetilde U}=N_{s'}$, de $\mathcal E_K$ (qui est extension d'un $\mathcal E_k$ par un facteur trivial \unsure{kummérien|kummerien} $T\simeq I_{s'}$, où $k$ est un corps de degré de \unsure{transcendance} $\unsure{\deg\operatorname{tr}K-1|\deg\operatorname{tr}n-1}$, lorsque $S$ est de type fini sur un corps de $\unsure{N}$, disons $\mathbf Q$) \uline{presque} un relèvement --- seul le facteur $I_{s'}\simeq T$ y \unsure{obstrue} encore : en \unsure{tous|tout} cas $N_{s'}\simeq\mathcal E_{\widetilde U}$ est plongé dans un $N_{\mathcal E_{X_s}}(L_i)$, et toute section de \unsure{cette} extension de $\unsure{E_{k(s)}|\mathcal E_{k(s)}}$ par $T$ donne donc une section de $\mathcal E_{X_s}$ sur $\mathcal E_{k(s)}$ --- mais de \uline{deuxième} espèce.

Une autre façon de dire ces choses est la suivante. Une section rationnelle $f$ de $\frac{U}{S}$ sur $S$, de domaine de définition l'ouvert $S'$ de $S$, \deleted{\unsure{peut}} peut s'interpréter en termes de la famille des pts $f(s)\in U_s$, où $s$ parcourt l'ens des pts de $S$
\margin{90}
\scanpage{92}
\unsure{-- ou} seulement, \unsure{sans} doute, l'ens des pts fermés (si $S$ est de type fini sur un corps, tel $\mathbf Q$ disons). En termes de sections d'extensions, on peut dire que la section $\unsure{\mathcal E_{f_\eta}|\mathcal E_{f_K}}$ de $\unsure{\mathcal E_{U_K}|\mathcal E_{U_\eta}}$ sur $\unsure{\mathcal E_K|\mathcal E_\eta}$ donne naissance à une famille de sections des $\mathcal E_{X_s}$ sur les $\mathcal E_{k(s)}$ pour $s\in S'$. On peut \deleted{\unsure{interpréter}} \inserted{\unsure{prolonger}} \unsure{ainsi|aussi} cette correspondance entre sections de $\unsure{\mathcal E_{U_\eta}|\mathcal E_{U_K}}$ sur $\unsure{\mathcal E_\eta|\mathcal E_K}$, et sections des $\mathcal E_{X_s}$ sur les $\mathcal E_{k(s)}$, en généralisant ce qu'on a vu dans le cas où $s$ est de codim $1$ : On introduit le normalisé $S'$ de $S$ dans $\overline K$, la fibre en $s$ s'identifie à l'ens des idéaux maximaux de la clôture \unsure{entière|intégrale} de $\underline{\mathcal O}_{S,s}$ dans $\overline K$, et si $N_{s'}$ désigne le \deleted{\unsure{normalisateur}} \inserted{stabilisateur} dans $\unsure{\mathcal E_{\overline K}|\mathcal E_{\overline\eta}}$ d'un $s'$ de $S'_s$, on a encore une structure d'extension
\[
(\unsure{20|29})\qquad
1\mathrel{\text{---}}I_{s'}\longrightarrow N_{s'}\longrightarrow \unsure{\mathcal E_{k(s)}|E_{k(s)}}\longrightarrow 1
\]
\margin{91}
\scanpage{93}
\mspage{46}
où $I_{s'}$ \deleted{est le groupe} \inserted{$\simeq\mathcal E_{\widetilde K}$}, $\widetilde{\widetilde K}$ étant le corps des fractions du hensélisé (\deleted{\illegible{2}} strict) de $\underline{\mathcal O}_{S,s}$. ($N_{s'}$ \deleted{est celui du corps de} \inserted{\unsure{s'identifie de|s'identifie à} $\unsure{=}$ à} $\mathcal E_{\widetilde K}$, où $\widetilde K$ correspond au hensélisé pas strict)

Ceci posé, on trouve comme précédemment, si le composé \deleted{de} $\mathcal E_K\mathrel{\unsure{\longrightarrow|-}}\mathcal E_{U_\eta}\mathrel{\unsure{\longrightarrow|-}}\mathcal E_U$ s'annule sur $I_{s'}$ \inserted{(a priori, il envoie $I_{s'}$ dans $\unsure{\Pi|T}$)} \deleted{\unsure{pour}} (ce qui signifie justement que $\unsure{s'|s}\in S'\unsure{\ldots}$) \deleted{que} \unsure{un hom|on trouve} $\mathcal E_{k(s)}\longrightarrow\mathcal E_U$, ou si on préfère $\mathcal E_{k(s)}\longrightarrow\mathcal E_{U_{\widetilde{\mathcal O}_s}}\xleftarrow{\sim}\mathcal E_{X_s}$, qui \deleted{peut} \inserted{donc} s'inter\unsure{prète|préter} comme une section de $\mathcal E_{X_s}$ sur $\mathcal E_{k(s)}$. Et que se passe-t-il pour $s\notin S'$ ? \deleted{\unsure{Alors}} On a vu que $D=S\setminus S'$ est purement de codim $1$ (c'est l'image réciproque du diviseur $T$ sur $X$ par la fonction $f$ de $X$ sur $S$, définie si $g\geq1$ ou si $S$ de dim $1\ldots$) en fait le $\unsure{\mu=\mu_\alpha|n=n_\alpha}$ de \unsure{tantôt}, relatif à
\margin{92}
\scanpage{94}
une composante $D_\alpha=\overline{\{s_\alpha\}}$ de ce diviseur, n'est autre que sa \uline{multiplicité} en $D_\alpha$.

Ceci rappelé, on se demande donc de ce qui se passe pour les $s\in\bigcup D_\alpha$, \deleted{On} \inserted{et} \deleted{voit que} regardons ceux qui sont $\in D_\alpha$, $\alpha$ fixé, i.e. qui spécialisent un $s_\alpha$ fixé. On peut dire que la situation définit une extension $\mathcal E_{U_{s_\alpha}}$ de $\mathcal E_{k(s_\alpha)}$ par $\pi$ --- déduite, comme tantôt pour $S$ et $K$, d'une courbe \unsure{relative} $U_{D_\alpha}$ sur $D_\alpha$ --- ou $U_{S_\alpha}$ sur un ouvert régulier convenable de $D_\alpha$, si on préfère --- en passant au pt générique. Cette fois-ci, pour la fibre générique, on n'a plus des \deleted{\unsure{sections}} pts. rat. sur $k(s_\alpha)$ de $\unsure{U_{s_\alpha}|U_{S_\alpha}}$ --- on n'a qu'un pt à l'infini $i_\alpha\in X_{s_\alpha}\setminus U_{s_\alpha}$, rat sur $\unsure{k_\alpha}$. Mais on dispose \inserted{maintenant} d'un $L_i\subset\pi$ relatif à cet $i_\alpha$, justement, de façon ess. canonique,
\margin{93}
\scanpage{95}
\mspage{47}
qui nous incite justement à regarder les classes de scindages de $\mathcal E_{\unsure{X_{s_\alpha}|X_{\overline s}}}$ sur $\mathcal E_{k(s)}$ relatifs à cet $L_i$ --- et notamment ceux relatifs au choix d'un uniformisant de $\mathcal O_{S,s_\alpha}$, i.e. d'une équation locale de $D_\alpha$ au vois. de son pt gén. $s_\alpha$. \deleted{\unsure{(plus gén.)}} \inserted{(qui intervient par le \unsure{biais|bout} qu'\unsure{elles définissent|elle définit} dans $\mathfrak m_{s_\alpha}/\mathfrak m_{s_\alpha}^{2}$)} Quant à ces sections-là, on voit comment elles se spécialisent, du moins c'est clair en les pts de $D_\alpha$ en lesquels la section rationnelle envisagée du faisceau conormal $\underline N_{D_\alpha,S}^{\unsure{\vee|*}}$ (si on suppose $D_\alpha$ régulier, \unsure{\&|et} donc $\underline N_{D_\alpha/S}$ inversible) reste régulière, et inversible.

On peut se proposer d'étudier \deleted{comment} \inserted{ce qui} se passe en les autres points --- i.e. la spécialisation des sections \inserted{de deuxième type} d'extensions de groupes profinis, en faisant une étude analogue à la précédente, pour cette question. Mais cela me semble \deleted{moins} \inserted{d'un} intérêt moins brûlant --- car cela n'est
\margin{94}
\scanpage{96}
plus une question pour des courbes relatives \uline{globales}, mais pour des ``éléments de courbe'' relatifs --- savoir, pour le complété formel d'une courbe relative $X/B$ le long d'une section $g$, avec le diviseur de cette section distingué dans la structure.

Je vais plutôt revenir à la situation où on part de une section de première espèce d'un $\mathcal E_{U_\eta}$ sur $\mathcal E_\eta$, \unsure{associée|associé} donc à un pt \unsure{rationnel|rationnelle} $f_\eta$ de $U_\eta$ sur $K$. On peut donc dire (indépendamment de la donnée du modèle $S$ de $K$ sur lequel $U_\eta$ pouvait déjà se définir en \deleted{\unsure{tant}} que schéma relatif élémentaire) que la dite section \deleted{de} $\mathcal E_{U_\eta}$ sur $\mathcal E_\eta$ \inserted{$(\unsure{29|23})\quad\unsure{\mathcal E f_\eta|f_\eta}=\psi:\mathcal E_\eta\longrightarrow\mathcal E_{U_\eta}$} définit, pour \unsure{tte|tt} section de $\mathcal E_K$ sur $\mathcal E_{\mathbf Q}=\Gamma_{\mathbf Q}$ ou sur un sous-groupe ouvert \inserted{$\Gamma'$} de celui-ci, une section
\margin{95}
\scanpage{97}
\mspage{48}
correspondante de $\mathcal E_{U_{\unsure{\eta|K}}}$ sur $\Gamma'$ --- d'où, en particulier, des actions de $\Gamma'$ sur $\pi$.

(Même chose d'ailleurs pour des sous-corps $k'$ quelconques de $K$, quand on regarde $\mathcal E_K$ comme étant sur $\mathcal E_{k'}$ : la section (29) donne une famille de sections de $\mathcal E_{U_{\unsure{\eta|K}}}$ sur des $\mathcal E_{k'}$, \inserted{ou des sous-groupes ouverts $\Gamma'$ de ceux-ci,} en la composant avec des sections de $\mathcal E_K$ sur des $\mathcal E_{k'}$ --- des sous-groupes $\Gamma'\subset\mathcal E_{k'}$. J'ai envie d'avoir une vue d'ensemble sur cette famille de sections, et de la mettre en relation avec la famille de sections des $\mathcal E_{X_s}$ sur les $\mathcal E_{k(s)}$, \inserted{définie par (29),} définissant des actions \unsure{effectives} (pas seulement ext) des $\mathcal E_{k(s)}$ sur $\pi$. Pour ces dernières, les $k(s)$ n'apparaissent pas comme des sous-corps de $K$, mais plutôt comme des $\underline{\mathcal O}_s/\mathfrak m_s$ --- des corps résiduels de sous-anneaux. Cependant, si on suppose $S$ de dim \deleted{\unsure{finie}} \inserted{\unsure{rel.|relative} 1} sur
\margin{96}
\scanpage{98}
\margin{On aura intérêt ici de prendre pour $S$ le modèle complet régulier can. de $K$ sur $k\ldots$}
un corps $k$, alors les $k(s)$ (pour $s\ne\eta$, pt gén.) apparaissent comme des extensions finies de ce $k$, et tout $s\in S$ définit une classe de conjugaison de sous-groupes $\unsure{I_{s'}|\Delta_{s'}}$ dans $\mathcal E_K$ (pour sa structure : \unsure{cette|la même} \inserted{\unsure{entière|extérieure}} relation sur $\mathcal E_k\ldots$), \inserted{et les $N_{s'}$ qui} \deleted{qui} sont justement les sous-groupes \deleted{normalisateurs} de \deleted{inertie} \inserted{décomposition} de tantôt, et on peut regarder les sections de \deleted{l'extension} \inserted{de $N_{s'}$ sur} des $\mathcal E_{k(s)}\underset{\inserted{\text{ouvert}}}{\subset}\mathcal E_k$ sur (sections donc ``de deuxième espèce'') et les composer avec la section donnée. Dans le cas où \deleted{$f$} \inserted{la section rat. $f$ de $U_K$ sur $S$} est \uline{définie} en $s$, il résulte aussitôt de ce qui précède que la section \inserted{partielle} obtenue de $\mathcal E_{U_K}$ au dessus de $\mathcal E_{k(s)}\hookrightarrow\mathcal E_{\unsure{K|k}}$ est aussi de première espèce, et \deleted{elle est} (en tant qu'op. \deleted{ext.} \inserted{effective} sur $\pi$) \deleted{\unsure{est}} n'est autre que celle qui est définie par le pt $f(s)\in U_s(k(s))$. Dans
\margin{97}
\scanpage{99}
\mspage{49}
le cas au contraire où \deleted{\unsure{il}} $f$ n'est pas définie en $s$, on trouve une \deleted{\unsure{section|seule}} action de $2^{\text{ième}}$ espèce, normalisant un $L_i$ \deleted{\unsure{cyclique}} \unsure{de|dans} $\pi$.

À vrai dire, on n'a regardé jusqu'à présent que ce qui se passe en des ``places'' $s$ de $K/k$ (corps des fonctions relatif d'une variable) où la \deleted{famille} \inserted{courbe} $\unsure{U_\eta/\eta|U_K/K}$ a ``bonne réduction'' --- c'était implicite, car on partait $=$ d'une $U$ sur $S$ --- et je voudrais examiner maintenant ce qui se passe dans le cas où on abandonne cette hypothèse : examiner quels types d'actions de $\mathcal E_{k(s)}\underset{\inserted{\text{ouvert}}}{\subset}\mathcal E_k$ on obtient par la construction précédente, en les $s\in S$ où on a ``mauvaise réduction''. Un cas type serait celui où $k=\mathbf Q$, $K$ corps des fonctions d'une courbe modulaire pour courbes elliptiques (\deleted{\unsure{considérées comme}} \inserted{pointées}), et $s$ place à l'infini de $K$, correspondant à une courbe elliptique
\margin{98}
\scanpage{100}
dégénérée. On peut passer même à l'hensélisé (non strict) de $\underline{\mathcal O}_s$, on \unsure{trouve|dans} la situation d'un trait hensélien $S$, à corps résiduel $k$ (ici $\mathbf Q$ --- ou une extension finie de $\mathbf Q$), d'une courbe $U_\eta$ de type $(g,\nu)$ (ici $(1,1)$) sur son pt générique $\eta$, on s'intéresse à l'extension $\mathcal E_{U_\eta}$ de $\mathcal E_\eta$ par $\pi$, où ici $\mathcal E_\eta$ s'inscrit encore dans
\[
(30)\qquad 1\longrightarrow T\longrightarrow\mathcal E_\eta\longrightarrow\mathcal E_k\longrightarrow1
\]
\deleted{On est ici en \unsure{car.} $0$} \inserted{en plus de}
\[
(31)\qquad 1-\pi-\mathcal E_{U_\eta}\longrightarrow\mathcal E_\eta\longrightarrow1
\]
Ici l'action extérieure de $T$ sur $\pi$ n'est sûrement pas triviale (je conjecture qu'elle l'est ssi il y a ``bonne réduction'') --- donc le groupe $\mathcal E_k$ n'opère pas lui $=$ extérieurement sur $\pi$. Mais tout scindage de (30)
\margin{mais je l'ai démontré, dans le contexte des VA, il y a belle lurette ! \unsure{Ça doit bien} être connu sur $H_1$ sûrement \unsure{!!}}
\margin{99}
\scanpage{101}
\mspage{50}
définit une extension de $\mathcal E_k$ par $\pi$, donc une action extérieure --- est-elle ``admissible'' ?, i.e. définie par une courbe algébrique ? Sans doute pas --- \uline{quelle} courbe serait-ce, si ce n'est la courbe ``réduite'' de type $(g,\nu)$ (inexistant !) ?
\margin{mais il y a tjrs une ``arithmétisation'' de $\pi$ \unsure{innocente} !}

Dès lors, si on a une section de (30) qu'on compose avec une de (31) --- donc une action \inserted{effective} de $\mathcal E_k$ sur $\pi$ --- qu'attendre d'une telle action ?

Sans m'y attendre, me voici à tâter un peu la structure à l'$\infty$ des variétés modulaires de Teichmüller !

S'il faut hasarder des conjectures dans cette situation \unsure{scabreuse}, ce pourrait-être celle-ci :

\uline{Conjecture-à-tout-hasard} : Les conditions suivantes sont équivalentes

\noindent a) $U_\eta$ a bonne réduction sur $S$

\noindent b) L'action extérieure de $T$ sur $\pi$ est triviale
\margin{100}
\scanpage{102}
Ce qui signifie aussi que $\forall$ \inserted{(ou un)} scindage de 31 --- p.ex. défini par un pt de $U_\eta$ rat./$\unsure{K|k}$ --- induit un hom $T\longrightarrow\pi$)

\noindent c) L'action de $T$ sur $\pi_{\mathrm{ab}}=H_1(U_{\bar\eta})$ est triviale.

\noindent d) Idem pour $H_1(U_{\bar\eta},\unsure{T_\ell|\mathbf Z_\ell})$, \unsure{si|où} $\ell$ est premier à la car. résiduelle. (\uline{NB} \deleted{\unsure{Pour car.} $p>0$} \inserted{$\ell=p>0$}, on doit pouvoir l'exprimer par un critère sur les groupes $p$-divisibles\ldots)

\noindent e) (En termes d'une section de (30), définissant une extension de $\mathcal E_k$ par $\pi$, et supposant maintenant $k$ de type fini sur $\mathbf Q$) Ladite extension est admissible (quitte à passer à un sous-groupe ouvert, \inserted{de $\mathcal E_k$,} au besoin), i.e. peut être définie par une courbe alg. sur \inserted{$k$, ou sur} une extension finie de $k$.

\noindent f) (En termes d'une section de (30) et d'une de (31), permettant de faire opérer $\mathcal E_k$ sur $\pi$) : Cette action \inserted{n'}est \inserted{pas trop} loufoque (à définir convenablement !)
\margin{101}
\scanpage{103}
\mspage{51}
On a évidemment $\mathrm{a)}\Rightarrow\mathrm{b)}\longrightarrow\mathrm{c)}\unsure{=|\Rightarrow}\mathrm{d)}$, et $d\Rightarrow a$ devrait être plus ou moins connu, ou du moins se prouver à l'aide de choses bien connues. Ces conditions impliquent e, mais je n'ai pas d'idée pour le moment comment voir que e) implique bonne réduction --- faudrait une compréhension (\unsure{fut-elle} heuristique) plus approfondie que celle que j'ai pour le moment.
\margin{Ça vaudrait la peine de regarder une \unsure{version} de ceci où $X_\eta=\bar U_\eta$ a bonne réduction --- mais la multisection $\overline{X_\eta-U_\eta}$ de $X$ sur $S$ n'est pas étale. Le cas du genre $0$ devrait être particulièrement aisé à regarder\ldots}

J'ai donc là un principe général (que je pourrais à l'occasion formuler avec la généralité qui lui convient) pour construire des actions \unsure{extérieures|\deleted{extérieures}} \inserted{Individuelles} de groupes $\mathcal E_k$ ($k$ \deleted{ext.} de type fini de $\mathbf Q$) sur des $\pi$ à \unsure{bouts}, qui \inserted{(sans doute)} ne soient géométriques, par \deleted{des} la considération de courbes de type $g,\nu$ ``se réduisant mal''. Mais je \unsure{n'arrive|reviens}
\margin{102}
\scanpage{104}
pas par cette voie à faire des actions \deleted{\unsure{par}} effectives, associées à \deleted{\unsure{une}} actions ext. déjà géométriques, qui ne soient elles-mêmes géométriques --- i.e. d'un des deux types $1^\circ$, $2^\circ$ envisagés au début de ce n°.
\margin{103}
\scanpage{105}
\mspage{52}
4. \uline{Sections \inserted{d'extensions,} et anneaux de valuation généraux.}

D'abord une révision de notations. Si $X$ est un schéma connexe, je \deleted{\unsure{dé}}note \deleted{$\mathcal E_X$}
\[
(1)\qquad \mathcal E_X=\pi_1(X)
\]
son groupe fond. profini en tant que groupe extérieur, et si $\tilde X$ est un rev. universel profini de $\mathcal E_X$, par
\[
(2)\qquad \mathcal E_X^{(\tilde X)}=\pi_1(X;\tilde X)=\operatorname{Aut}_X(\tilde X)
\]
son groupe fondamental précisé --- qui est un groupe profini. Si $X=\operatorname{Spec}(A)$, où $A$ est un anneau (le plus souvent un corps) je note $\mathcal E_A$, et $\mathcal E_A^{\tilde A}$ \inserted{$=\mathcal E_X^{\unsure{\tilde X|X}}$} si $\tilde A$ est une $A$-algèbre telle que $\operatorname{Spec}(\tilde A)=\tilde X$ soit un rev. universel de $X$ (ce qui le détermine à isom non unique près). Bien entendu, si $\unsure{\xi|\zeta}$ est un ``point géométrique'' de $X$, on note
\[
(3)\qquad \mathcal E_X^{\unsure{\xi|\zeta}}=\pi_1(X,\unsure{\xi|\zeta})=\mathcal E_X^{\tilde X(\unsure{\xi|\zeta})},
\]
où $\tilde X(\unsure{\xi|\zeta})$ est le rev. universel de $X$ pointé en $\unsure{\xi|\zeta}$. Le choix
\margin{104}
\scanpage{106}
de $\unsure{\xi|\zeta}$ correspond au choix d'un $x\in X$ et d'une ext. sép. close \deleted{$\bar K$} \inserted{$\Omega$} de $k(x)$ (pas néc. \deleted{\unsure{une}} clôture alg. de $k(x)$) et on note alors aussi $\mathcal E_X^\Omega$ au lieu de $\mathcal E_X^{\unsure{\xi|\zeta}}$ ($\Omega$ sous-entendu en tant qu'extension de $k(x)$, donc avec sa structure de $k(x)$ algèbre). On aura \inserted{un isom can.}
\[
(4)\qquad \mathcal E_{\unsure{\alpha|X}}^\Omega\simeq\mathcal E_X^{\overline{k(x)}},
\]
où $\overline{k(x)}$ est la clôture alg. de $k(x)$ dans $\Omega$. Bien sûr, si $X=\operatorname{Spec}(A)$, on notera aussi $\mathcal E_A^\Omega$ --- notation surtout utilisée pour $\mathcal E_K^{\bar K}$, $K$ un corps, $\bar K$ une clôture alg. ou sép. de $K$.

Si $X$ est un $Y$-schéma, \inserted{$Y$ connexe,} on notera
\[
(5)\qquad \mathcal E(f):\mathcal E_X\mathrel{\text{---}}\mathcal E_Y\qquad\text{où }f:X\longrightarrow Y
\]
($\mathcal E_X$ est un foncteur en $X$
\[
(\unsure{\text{Schémas}|\text{Sch conn}})\mathrel{\text{---}}\unsure{\text{grpes ext}}
\]
\margin{105}
\scanpage{107}
\mspage{53}

qui se précise par l'hom \unsure{effectif|fonctel} de groupes profinis
\[
(6)\qquad \mathcal E_X^{\widetilde X}\longrightarrow\mathcal E_Y^{\widetilde Y}
\]
où $\widetilde Y$ est le rev. universel de $Y$ défini par $\widetilde X\longrightarrow Y$ ($\widetilde X$ jouant le rôle du point --- on devrait pouvoir écrire en fait $\mathcal E_Y^{\widetilde X}$, plus gén. $\mathcal E_Y^Z$ chaque fois qu'on a un $Y$-schéma $Z$ $1$-connexe, jouant le rôle de ``foncteur fibre'' pour le topos $B_{\unsure{\pi(X)}}$ des rev étales de $Y$\unsure{.})

\deleted{\unsure{que} $\mathcal E_X^{\widetilde X}$}

On peut dire que
\[
(7)\qquad (X,\widetilde X)\longmapsto\mathcal E_X^{\widetilde X}
\]
est un foncteur, de la catégorie des schémas $0$-connexes $X$ \uline{munis} d'un rev. universel (ou seulement d'un $Z$ $1$-connexe s'envoyant dans $X$) vers celle des groupes profinis. Ceci s'applique en particulier en regardant la \inserted{sous-}cat. des $(X,\unsure{\xi|\zeta})$ munis d'un $p^t$ géom ---

\margin{106}
\scanpage{108}

on a donc
\[
(8)\qquad \mathcal E_X^{\unsure{\xi|\zeta}}\longrightarrow\mathcal E_Y^\eta
\]
si on a un hom de schémas ``\unsure{géom\textsuperscript{t}.} \unsure{ponctués|pointés}'' $(X,\unsure{\xi|\zeta})\longrightarrow(Y,\eta)$. On notera que \unsure{tte|tt} ponctuation géom. de $X$ en un $x\in X$ --- i.e. une ext. sep. close $\Omega$ de $k(x)$ --- définit \inserted{\unsure{can\textsuperscript{t}}} une ponctuation \inserted{géom.} aussi de $Y$ en $y=f(x)$, \deleted{par} par composition --- $\Omega$ ``est'' une extension de $k(y)$ puisque $\unsure{k(x)}\hookleftarrow k(y)$ --- et l'hom (8) s'identifie aussi à
\[
(9)\qquad \mathcal E_X^{\overline{k(x)}}\longrightarrow\mathcal E_Y^{\overline{k(y)}},
\]
où $\overline{k(x)}$, $\overline{k(y)}$ sont les clôtures séparables \uline{dans $\Omega$}.

On posera
\[
(10)\qquad \mathcal E_{X/Y}^{\widetilde X}=\operatorname{Ker}\bigl(\mathcal E_X^{\widetilde X}\longrightarrow\mathcal E_Y^{\widetilde X}=\mathcal E_Y^{\widetilde Y}\bigr)
\]

\margin{107}
\scanpage{109}
\mspage{54}

(C'est un foncteur \unsure{p.r.} au triple $\widetilde X\mathrel{\text{---}}X\mathrel{\text{---}}Y$ avec $X,Y$ $0$-connexes, $\widetilde X$ un rev. universel; plus \unsure{généralt}, si $\unsure{T|Z}\unsure{\mathrel{\text{---}}|\longrightarrow}X$ avec $\unsure{T|Z}$ $1$-connexe, on pose
\[
(11)\qquad \unsure{\mathcal E}_{X/Y}^{\unsure{T|Z}}=\operatorname{Ker}\bigl(\mathcal E_X^{\unsure{T|Z}}\unsure{\mathrel{\text{---}}|\longrightarrow}\mathcal E_Y^{\unsure{T|Z}}\bigr)
\]
on a un foncteur en \inserted{\unsure{le triple}} $(\unsure{T|Z}\unsure{\mathrel{\text{---}}|\longrightarrow}X\longrightarrow Y)$ variable. En particulier: $\mathcal E_{X/Y}^{\unsure{\xi|\zeta}}$, $\unsure{\xi|\zeta}$ un $p^t$ géom. de $X$, $\mathcal E_{X/Y}^{\Omega}$, $\mathcal E_{X/A}^{\widetilde X}$ (si $Y=\operatorname{Spec} A$), $\mathcal E_{B/A}^{\widetilde B}$ $\ldots$

\margin{\uline{NB} $\mathcal E_{X/Y}^T$, pour $T$ variable, $X/Y$ fixé, est défini à isom non unique \unsure{!} près, et se note $\mathcal E_{X/Y}$ \unsure{simplem\textsuperscript{t}.}}

Un cas intéressant est celui où on dispose d'une ``suite exacte d'homotopie universelle'' \deleted{de} (en dim $1$) pour $X\longrightarrow Y$, alors la donnée (pour $X=Y$ donné) de $T\unsure{\mathrel{\text{---}}|\longrightarrow}X$, (avec $T$ $1$-connexe) peut s'interpréter par la donnée \unsure{du} \unsure{composé}

\margin{Car $\mathcal E_X^T\longrightarrow\mathcal E_Y^T$ est épi.}
\margin{108}
\scanpage{110}

\[
T\longrightarrow Y
\]
\uline{et} d'un relèvement en $T\longrightarrow X$, ou ce qui revient au $=$, d'une section de $X_T=X\times_YT$ sur $T$. Ceci posé;

\diagram{\begin{tikzcd}
X\arrow[d] & X_T\arrow[l]\arrow[d] & \\
Y & T\arrow[l] & T\arrow[ul]\arrow[l,equal]
\end{tikzcd}}

on aura un isom \unsure{can|can\textsuperscript{t}} (sous l'hyp. de ``suite exacte d'homotopie'' faite)
\[
(12)\qquad \mathcal E_{X/Y}^T\simeq\pi_1(X_T;T)\simeq\mathcal E_{X_T}^T
\]
et on aura
\[
(13)\qquad 1\longrightarrow\mathcal E_{X/Y}^T\unsure{\mathrel{\text{---}}|\longrightarrow}\mathcal E_X^T\unsure{\mathrel{\text{---}}|\longrightarrow}\mathcal E_Y^T\unsure{\mathrel{\text{---}}|\longrightarrow}1
\]
(Cette hypothèse sera tjrs satisfaite si $Y=\operatorname{Spec}K$, $K$ un corps, si $X$ est géom. $0$-connexe sur $K$)

Plus généralement, si on a une suite exacte d'homotopie universelle, mais pas néc. avec $\mathcal E_X^T\unsure{\mathrel{\text{---}}|\longrightarrow}\mathcal E_Y^T$ surj,

\margin{109}
\scanpage{111}
\mspage{55}

\margin{Sous l'hyp. ``suite exacte d'homotopie'', mais avec fibres pas néc. \unsure{conn.}}

on aura une factorisation can. de $X\longrightarrow Y$ en
\[
(14)\qquad X\longrightarrow Y'\longrightarrow Y
\]
avec $Y'$ étale fini ou pro-étale fini sur $Y$ et $\mathcal E_X^{\unsure{\prime}}\longrightarrow\mathcal E_{Y'}$ est maintenant épi, \deleted{\unsure{donc}} \unsure{cette|avec} suite exacte univ. d'homotopie bien sûr. On aura des iso can
\[
(15)\qquad \mathcal E_{X/Y}^T\xrightarrow{\sim}\mathcal E_{X/Y'}^T
\]
qui peut donc se calculer comme
\[
(16)\qquad \mathcal E_{X\times_{Y'}T}^T\simeq\mathcal E_{X/Y}^T
\]
qu'on peut noter $\mathcal E_{X_T}^T$, mais en faisant attention que maintenant $X_T$ désigne non plus $X\times_YT$ (qui va être \unsure{disconnexe} \inserted{si $Y'\longrightarrow Y$ pas iso}) mais $X\times_{Y'}T$.

Bien sûr, à iso. can. près on peut \deleted{\unsure{en}} remplacer, pour le calcul de $\mathcal E_{X_T}^T$, $T$ par $\widetilde Y$ pour ce qui est du $X_T$, et par $\widetilde X$ pour ce qui est de \unsure{l'ancrage|l'origine} dans $X_{\widetilde Y}$,

\margin{110}
\scanpage{112}

où $\widetilde X,\widetilde Y$ sont les rev. universels de $X,Y$ déduits de \inserted{$T\longrightarrow X$ et de $T\longrightarrow Y$,} de telle façon que l'on obtient
\[
(17)\qquad \mathcal E_{X/Y}^T\simeq\mathcal E_{X/Y}^{\widetilde X}\simeq\operatorname{Ker}\bigl(\mathcal E_X^{\widetilde X}\unsure{\mathrel{\text{---}}|\longrightarrow}\mathcal E_Y^{\widetilde Y}\bigr)
\]
(sans hyp. de suite exacte)

et
\[
\begin{aligned}
(18)\qquad \mathcal E_{X/Y}^T\simeq{}&\deleted{\operatorname{Ker}\bigl(\unsure{\mathcal E_{X_{\widetilde Y}}^{\widetilde X}}\unsure{\mathrel{\text{---}}|\longrightarrow}\unsure{\mathcal E_{\widetilde Y}^{\widetilde Y}}\bigr)}\\
&\inserted{\mathcal E_{X_{\widetilde Y}}^Z\unsure{\simeq}\mathcal E_{X_{\widetilde Y}}^{\widetilde X}}
\end{aligned}
\]
dans le cas général.

\diagram{\begin{tikzcd}
X\arrow[d] & X_{\widetilde Y}\arrow[l]\arrow[d] & \widetilde X\arrow[l] & Z\arrow[l]\\
Y'\arrow[d] & \widetilde Y\arrow[l] & & \\
Y & & &
\end{tikzcd}}

Une autre interprétation de $\mathcal E_{X/Y}^T$ s'obtient en prenant $T$ réduit à un $p^t$ géométrique \inserted{$\unsure{\xi|\zeta}_X$}, qui définit donc un $p^t$ géom. $\unsure{\xi|\zeta}_Y$ de $Y$, une fibre géom. $X_{\unsure{\xi|\zeta}_Y}$ de $X$ sur $Y$ en $\unsure{\xi|\zeta}_Y$, $\unsure{\xi|\zeta}_X$ s'interprétant comme un $p^t$ géom. de cette fibre, et on aura
\[
(19)\qquad \mathcal E_{X/Y}^{\unsure{\xi|\zeta}_X}\simeq\mathcal E_{X_{(\unsure{\xi|\zeta}_Y)}}^{\unsure{\xi|\zeta}_X}\simeq\pi_1(X_{(\unsure{\xi|\zeta}_Y)},\unsure{\xi|\zeta}_X)
\]
--- i.e. $\mathcal E_{X/Y}$ s'interprète comme groupe fondamental des ``fibres géom'' de $X/Y$.

\margin{111}
\scanpage{113}
\mspage{56}

Le lien avec la \deleted{définition}\inserted{description} (11) est évidemment que, si $\unsure{\xi|\zeta}_X$ provient d'un $p^t$ géom $\unsure{\xi|\zeta}_T$ de $T$, alors on aura des iso can

(20)\quad
\diagram{\begin{tikzcd}[column sep=large,row sep=large]
\mathcal E_{X/Y}^{\unsure{\xi|\zeta}_T}\arrow[r,phantom,"\simeq"]\arrow[d,phantom,"\rotatebox{90}{$\simeq$}" description] & \mathcal E_{X/Y}^T\\
\pi_1(X_{\unsure{\xi|\zeta}_Y},\unsure{\xi|\zeta}_T) &
\end{tikzcd}}

\deleted{\unsure{donc}} --- les $\pi_1(X_{\unsure{\xi|\zeta}_Y},\unsure{\xi|\zeta}_T)$ des fibres géom, relatifs aux divers pts géom de $\mathcal E_{X/Y}^T$, sont canoniquement iso entre eux.

( Il faut bien faire attention que $\mathcal E_{X/Y}$ \inserted{($\unsure{Z|T}$ non précisé)} est une classe d'iso. de groupes profinis sans plus, \uline{pas} un groupe extérieur. L'arbitraire dans le choix d'un iso entre $\mathcal E_{X/Y}^T$ et $\mathcal E_{X/Y}^{T'}$ provient du choix

\margin{112}
\scanpage{114}

d'un él. de $\mathcal E_{X\deleted{\unsure{/Y}}}^T$, opérant sur le sous-groupe $\mathcal E_{X/Y}^T$ par autom. intérieurs.

Les notations \inserted{$\mathcal E_X$,} $\mathcal E_{X/Y}$ évoquent intuitivement un $p^t$ de vue qui regarde $\mathcal E_{X\deleted{\unsure{/Y}}}^T$ comme une \uline{extension} de $\mathcal E_Y^T$ (ou d'un ss groupe convenable \unsure{seulement}) par $\mathcal E_{X/Y}^T$, qui sera souvent, pour les besoins de la cause, noté $\pi$ --- il joue le rôle d'un groupe fondamental (géométrique), tandis que $\mathcal E_Y^T$ joue le rôle d'un groupe de Galois --- ou d'un groupe profini mixte.

\margin{et j'hésite encore entre les \uline{notations} $\mathcal E_{X/Y}^T$ et $\pi_{X/Y}^T$; que je vais peut-être employer concurremment.}

Le cas le plus important pour nous sera celui où $Y=\operatorname{Spec}K$, $K$ un corps. \deleted{\unsure{Assez} souvent} Dans ce cas, $\mathcal E_Y\simeq\mathcal E_Y^{\widetilde Y}$

\margin{113}
\scanpage{115}
\mspage{57}

s'identifie bel et bien au groupe de Galois $\operatorname{Gal}(\overline K/K)$ de \deleted{$\unsure{\mathcal E_Y}$} $\overline K/K$, $\overline K$ étant la clôture sép. de $K$ telle que $\operatorname{Spec}\unsure{\widetilde K|\overline K}=\widetilde Y$. Souvent, on notera $\Gamma$, \inserted{ou $\Gamma_K$, $\Gamma_K^{\overline K}$,} au lieu de $\mathcal E_Y$ --- surtout si $K$ est alg. sur le corps premier, et qu'il n'a donc plus de ``partie géométrique'' à distinguer d'une ``partie arithmétique''$\ldots$

Soit \inserted{\unsure{un}} $K$ \deleted{\unsure{tel}} corps (qui pourrait être alg. clos), $L$ une ext. de type fini de $K$, $X$ un ``modèle'' propre régulier de $L$. Alors $\mathcal E_X^{\overline L}$ s'identifie à un quotient de $\mathcal E_L^{\overline L}$, qui \uline{ne \deleted{\unsure{p}} dépend pas du modèle $X$ choisi}, comme il est bien connu

\margin{114}
\scanpage{116}

c'est la partie ``universellement géom'' de $\mathcal E_L^{\overline L}$, \deleted{noté $\mathcal E_{L/K}^{\unsure{g},\overline L}$} \deleted{qui \unsure{couvre} les \unsure{hom|bons}} \inserted{(relativement à $K$),} qui classifie les schémas (finis) étales sur $L$ qui sont ``non ramifiés'' sur \uline{tout} modèle \deleted{\unsure{étal}} régulier (propre ou non) de $L/K$.

\margin{noté $\mathcal E_L^{\unsure{g},\overline L}$ ou $\mathcal E_L^{\unsure{g(K)},\overline L}$}

Si $U$ est un modèle quelconque; il se plonge dans un $X$, et on a des hom surjectifs
\deleted{\unsure{q}}
\[
(21)\qquad \mathcal E_L^{\overline L}\longrightarrow\mathcal E_U^{\overline L}\longrightarrow\mathcal E_X^{\overline L}=\mathcal E_L^{\unsure{g},\overline L}
\]

On aura \deleted{\unsure{alors}}
\[
(22)\qquad U=X\setminus Z
\]
$Z$ partie ferm. de $X$ --- soient $D_i$ ses composantes de codim $1$, $x_i$ les pts gén. de $D_i$. On a des groupes de décomposition et d'inertie $N_{x_i'}$ et $I_{x_i'}$ \inserted{de $\mathcal E_U^{\overline L}$}, (relatifs au choix d'un $x_i'$ au dessus de $x_i$ dans le normalisé $X'$ dans $\overline L$), et on trouve

\margin{115}
\scanpage{117}
\mspage{58}
\[
(23)\qquad
\mathcal E_X^{\deleted{\unsure{K}}\inserted{\overline L}}
=\mathcal E_L^{\unsure{g|G}\overline L}
\simeq\mathcal E_U^{\overline L}/\begin{gathered}\text{ss-groupe fermé}\\\text{inv. engendré}\\\text{par les }\unsure{L_{x_i'}|I_{x_i'}}\\(i\in\unsure{I|\overline I})\end{gathered}
\]
Notons que $\mathcal E_L^{\overline L}$ est la limite projective des $\mathcal E_U^{\overline L}$, pour des modèles réguliers (\unsure{affines si on y tient|affines non vides}) variables
\[
(24)\qquad\mathcal E_L^{\overline L}\simeq\lim_{\longleftarrow\,U}\mathcal E_U^{\overline L}
\]
et de \unsure{même}, bien sûr
\[
(25)\qquad\pi_{L/K}^{\overline L}\simeq\lim_{\longleftarrow\,U}\pi_{U/K}^{\overline L}.
\]
Comment récupérer les $\mathcal E_U^{\overline L}$ \unsure{à|à partir de} $\mathcal E_L^{\overline L}$ ? Pour \uline{\unsure{tt}} sous-anneau \inserted{\unsure{d-vr} $V$} de \unsure{ess.|en.} de type \unsure{fini} sur $K$ (correspond à un ``germe de modèle'' de $L/K$, au voisinage d'un pt de codim $1$ i.e. du pt gén. d'un diviseur), considérons le point fermé $x$ de $S_x=\operatorname{Spec}(V)$,
\margin{Dans les considérations qui suivent, $K$ (et \unsure{$X,U$}) ne jouent provisoirement aucun rôle --- (Mais on pourrait ``plus gén.'' \unsure{remplacer} \unsure{groupes|les groupes} \unsure{de décomp.|par des} \unsure{d'un des|l'un des} ``$U$'' \deleted{\unsure{de dimension relative}} \uline{\unsure{mais seulement}} \uline{\unsure{épi-s-comp|épi-scop}} (la ``\unsure{généralisation|géométrie}'' étant \unsure{triviale|prise}, $\unsure{\mathcal E_L\longrightarrow\mathcal E_U|\mathcal E_L^{\overline L}\longrightarrow\mathcal E_\eta}$))}
\margin{NB on s'intéresse en termes ``\unsure{absolus}'' \unsure{que pour des schémas} \uline{\unsure{épais}} (\unsure{quasi-proj.} ?) \uuline{réguliers}}
\margin{116}
\scanpage{118}
et un $x'$ au dessus de $x$ dans le normalisé de $S_x$ dans $\overline L$, d'où un groupe de décomposition $N_{x'}$, et un groupe d'inertie $\unsure{L_{x'}|I_{x'}}$, avec
\[
(26)\qquad N_{x'}/I_{x'}\simeq k(x')
\]
(ext. de type fini de $K$, de degré de tr. $n-1$, si $n=\deg\operatorname{tr}L/K$), et
\[
(27)\qquad I_{x'}\simeq\unsure{T_0|T_\infty}(\overline L)\simeq T_\infty(\unsure{\widetilde K|\overline K})\;-
\]
(par Kummer). \deleted{Notons \unsure{bien, on a}} \inserted{Ces constructions} ne dépendent du choix d'aucun \deleted{de} modèle ($U$, ou $X\ldots$). Notons que
\[
(28)\qquad I_{\unsure{x'|x}}\subset\pi_{L/K}^{\overline L}.
\]
En fait, on peut dire, posant
\[
(29)\qquad\left\{\begin{aligned}
\widetilde S_x&=\operatorname{Spec}\widetilde V,\quad\widetilde V\text{ hensélisé (pas strict) de }V\\
\widetilde U_x&=\widetilde S_x\setminus\{x\}\simeq\operatorname{Spec}\widetilde L_x\quad(\text{où }\widetilde L_x\text{ est le corps des \unsure{fractions} de }\widetilde V)
\end{aligned}\right.
\]
que l'on a un hom \inserted{ext} $\mathcal E_{\widetilde U_x}\longrightarrow\mathcal E_L$, \deleted{\unsure{alors}} dont le choix ``effectif'' dépend de
\margin{$I_{x'}$ \unsure{joue} le rôle d'un ``groupe bout''}
\margin{si Car. $k(x)$ est nulle --- \unsure{sinon quotient du seul habituel}\ldots}
\margin{117}
\scanpage{119}
\mspage{59}
celui d'un rev. universel \deleted{\unsure{(ou}} \inserted{ou encore} d'un pt géom $\unsure{\xi|\zeta}$ de $\unsure{\widetilde K_x|\widetilde U_x}$ --- i.e. \unsure{d'une|de une} clôture algébrique de $\widetilde K_x$ --- \uline{et} d'une \deleted{\unsure{isom}} ``chemin de \unsure{domaines|connexion}'' entre $\unsure{\xi|\zeta}_{\widetilde U_x}\unsure{\mathrel{\text{---}}|\longrightarrow}\unsure{\widetilde K_x|\widetilde U_x}\longrightarrow\operatorname{Spec}(L)$ et $\unsure{\xi|\zeta}_L$ défini par $\unsure{\widetilde L|\overline L}$ --- ou ce qui revient au \unsure{=|même}, \inserted{à isom près des \unsure{revts}} d'un $\unsure{K|L}$-hom de \deleted{$\unsure{K}$} $\widetilde L_x$ dans $\overline L$ --- et ceci à son tour équivaut à la donnée d'un $x'$, \deleted{\uline{et} \unsure{d'un is}} comme dessus. Ceci posé, on a des isom can.
\[
(30)\quad\left\{\begin{aligned}
N_{x'}&\simeq\mathcal E_{\widetilde U_x}^{\overline L}\\
I_{x'}&\simeq\text{Noyau de l'hom canonique }\inserted{\text{surjectif}}\\
&\mathcal E_{\widetilde U_x}^{\overline L}\longrightarrow\unsure{\mathcal E_{\overline{k(x)}}|\mathcal E_{k(x)}^{\overline{k(x)}}}
\end{aligned}\right.
\]
(où la clôture alg $\overline{k(x)}$ de $k(x)$ est induite par la clôture alg. $\overline L$ de $\widetilde K_x$)
\margin{Cf théorie galoisienne des traits henséliens. Ici $I_{x'}$ joue le rôle d'un $\pi_{\widetilde U_x/k(x)}^{\overline L}$ ; en fait il est $=\pi_{\widetilde U_x/\widetilde S_x}^{\overline L}$ !}
(\unsure{Alors} (27) est déduit par théorie de Kummer)
\deleted{$I_{x'}\simeq$}
\deleted{(\unsure{27|31})}

\margin{$\Vert$ ?}
Je présume que (dans la situation actuelle, ou même, avec un corps de base $K$) on a $N_{x'}=\operatorname{Norm}_{\mathcal E_L}(I_{x'})$, \inserted{(a priori, on a seulement $\subset$)} ce qui est lié à l'idée qu'on doit avoir $I_{x'}=\operatorname{Norm}_{\pi_{L/K}}(I_{x'})$
\margin{118}
\scanpage{120}
--- à tirer au clair à l'occasion !

Revenant au modèle $U$, fournissant un quotient $\mathcal E_U^{\overline L}$ de $\mathcal E_L^{\overline L}$, d'où $N_{x'}\longrightarrow\mathcal E_U^{\overline L}$, induisant $I_{x'}\longrightarrow\mathcal E_U^{\overline L}$,
\deleted{il est clair que si}
qui n'est autre que l'hom \deleted{\unsure{can|$\longleftarrow$}} $\mathcal E_{\widetilde U_x}\unsure{\mathrel{\text{---}}|\longrightarrow}\mathcal E_U$ associé à $\widetilde U_x\longrightarrow U$. \uline{Si} cet hom factorise par $\widetilde S_x$, il est immédiat que l'hom induit $I_{x'}\longrightarrow\mathcal E_U^{\overline L}$ est nul, puisque \inserted{$I_{x'}$} est le noyau de $\mathcal E_{\widetilde U_x}\longrightarrow\mathcal E_{\widetilde S_x}$, et $\mathcal E_{\widetilde U_x}\unsure{\mathrel{\text{---}}|\longrightarrow}\mathcal E_U$ se factorise par cet hom.
\margin{i.e. \uline{si} $V$ a un ``centre sur $U$'' i.e. domine un anneau local de $U$\ldots}
[ \uline{A t-on une réciproque ?} regarder ! ]
\margin{Ça devrait être le cas pour $U$ ``élémentaire'' \unsure{à|et} facteurs \unsure{pas|par} de \uline{\unsure{types $(0,0)$}} \uuline{\unsure{ni|ou}} \uline{$(0,1)$\ldots} \unsure{lié à ceci} \unsure{enfin à l'existence} d'un $K$}

Quand on dispose d'un modèle \uline{projectif} dont $U$ est un ouvert, on peut dire (par \unsure{définition valuative|critère valuatif} de la propreté) que $V$ a nécessairement un centre sur $X$, soit $a$, \deleted{et le centre} \inserted{(\uline{NB} $a$ pas} néc. de codim.1\ldots) et la condition envisagée tantôt (assurant que $I_{x'}\unsure{\mathrel{\text{---}}|\longrightarrow}\mathcal E_U^{\overline L}$ est nul) est que $a\in U$.
\margin{119}
\scanpage{121}
\mspage{60}
Ceci posé, $\mathcal E_U^{\overline L}$ se récupère à partir de $\mathcal E_L^{\overline L}$, comme quotient de ce dernier, en prenant \uline{tous} les $V$ de $L$ ayant un centre sur $U$ (il suffit \unsure{=|même} de prendre les $V=\underline{O}_{U,x}$, où $x$ est un pt de codim 1 dans $U$), et en divisant par le ss-groupe fermé invariant engendré par les $I_{x'}$ correspondants.

On peut regarder l'espace loc. annelé, $\varprojlim$ des modèles \inserted{(réguliers)} propres de $L$ sur $K$
\[
(31)\qquad\mathfrak X_{L/K}\simeq\varprojlim_{\substack{X\text{ modèle}\\\text{propre de }L/K}}X
\]
\margin{\uline{NB} Cette définition ne dépend pas de l'hypothèse de résolution des sing --- car on peut prendre des modèles normaux --- pas néc. réguliers --- (Ça revient au \unsure{=|même}, si on dispose de la résolution\ldots)}
--- ce n'est \uline{pas} un schéma ! \inserted{(si $n\geq2$)}. \deleted{\unsure{L'on}} \inserted{Si}
$x\in\mathfrak X$, \deleted{\unsure{l'anneau de}} \inserted{$k(x)$} est une ext. de $K$ de degré de transcendance $d\leq n$ --- on l'appelle \uline{dimension} de $x$, \inserted{et $n-d$ est la \uline{codim.}} Il y a un pt et un seul de codim 0 --- correspondant aux pts gén. des $X$. Les points de codim 1
\margin{120}
\scanpage{122}
correspondent exactement (via $x\longmapsto\underline{O}_{\mathfrak X,x}$) aux jauges dans $L$ qui sont \unsure{ess.|ext.} de type fini sur $K$, de corps des fractions $L$ --- i.e. aux $V$ précédents.

Sauf erreur, ces deux types de pts de $\mathfrak X$ (qui épuisent \deleted{tout} si $n=1$, \inserted{qui est le seul cas, avec $n=0$, où $\mathfrak X$ est un schéma !}) sont les seuls pour lesquels $\underline{O}_{\mathfrak X,x}$ soit noethérien. Cependant, on a

\margin{$\big|$}
\uline{Proposition} Pour $\forall x\in\mathfrak X$, $\underline{O}_{\mathfrak X,x}$ est un anneau de valuation dans $L$, de corps des fractions $L$, et contenant $K$. L'application $x\longmapsto\underline{O}_{\mathfrak X,x}$ est une bijection de $\mathfrak X$ sur l'ens de ces anneaux dans $L$.
\margin{On pourrait partir de un schéma excellent $X$, et construire un $\mathfrak X$ relatif au dessus\ldots}

\uline{Dém} a) Soit $f\in L\setminus\underline{O}_{\mathfrak X,x}$, il faut prouver $\frac1f\in\underline{O}_{\mathfrak X,x}$. Or $f=g/h$, \inserted{$g,h$}$\in\underline{O}_{X,a}$, où $X$ modèle \inserted{régulier} (propre), $a$ un pt, $f,g$ \deleted{premiers entre eux} --- utilisant la résolution, \deleted{\unsure{on} \uline{\unsure{sépare}} \unsure{les diviseurs}} \inserted{OPS $\exists$ \unsure{système} régulier} de paramètres, avec $f,g$ des monômes
\margin{121}
\scanpage{123}
\mspage{61}
en ces paramètres, et en se ramenant à un $f,g$ premiers entre eux, \unsure{on est|on voit} --- faisant encore éclater $a$ --- que les diviseurs de $f,g$ \deleted{\unsure{sont|se}} \unsure{séparent|séparés}\ldots

b) Un anneau de val $V$ \deleted{qui} pour l'extension $L/K$ domine un pt bien déterminé de chaque $X$, d'où un $x\in\mathfrak X=\varprojlim X$, et $V$ domine $\underline{O}_{\mathfrak X,x}$, donc lui est égal puisque ce dernier est un anneau de valuation.

\uline{NB} Tout \inserted{modèle $U$ de $L/K$ définit un} \uline{ouvert} de $\mathfrak X$ \deleted{\unsure{correspondant}} \deleted{\unsure{évidemment}} \inserted{(en le plongeant} dans un modèle propre $X$, c'est clair) et \deleted{et} les modèles correspondent à \uline{certains} ouverts de $\mathfrak X$ (\uline{pas} stables par réunions finies).

Pour \unsure{tt|$\forall$} anneau \unsure{loc\textsuperscript{l}|local} \inserted{\unsure{norm\textsuperscript{l}|normal}} $O\subset L$ de corps des fractions $L$, on peut regarder $\operatorname{Spec}\underline{O}=S$, son pt fermé $\unsure{x|s}$, le \unsure{normali-}
\margin{122}
\scanpage{124}
$O'$ de $O$ dans $\overline L$ i.e. $S'$ de $S$ dans $L$, et pour $x'$ au dessus de $x$, un groupe de décomposition $N_{x'}$, d'inertie $I_{x'}$, ayant les interprétations qu'on a signalé plus haut dans le cas particulier d'un $O$ qui est un anneau jauge. L'idée, c'est que prenant $O$ de valuation (pas néc. une jauge) $\supset K$, on trouve des groupes d'inertie à structure simple --- des $T_\infty^\alpha$ essentiellement, où $\alpha$ est la ``codim'' du $x\in\mathfrak X$ correspondant --- et si $\unsure{a|x}$ est de dim 0, on \unsure{aur.|aura} une extension de $\mathcal E_{k(x)}^{\overline{k(\unsure{a|x})}}$ --- ss-groupe ouvert de $\mathcal E_K$, puisque $k(\unsure{a|x})$ est \inserted{fini} alg. sur $K$ --- \unsure{par} $T^\alpha$, qu'on aura une chance de pouvoir scinder sans trop de mal. Ceci donne-
\margin{$T_\infty=T_\infty(\overline L)$}
\margin{on le supposera du moins ici}
\margin{123}
\scanpage{125}
\mspage{62}
rait un principe général pour construire des germes de sections de $\mathcal E_L$ sur $\mathcal E_K$ --- à l'aide des pts de dim 0 de la ``variété de Riemann'' $\mathfrak X$ de $L$ sur $K$. Composant, pour une sous-extension $L_1^{\unsure{k|*}}$ de $K$ de degré de tr. 1, avec l'hom $\mathcal E_L\longrightarrow\mathcal E_{L_1}$, on a donc aussi un moyen de construire des \inserted{germes de} sections de $\mathcal E_{L_1}$ sur $\mathcal E_K$, et le fait intéressant sera de voir si on \inserted{en} trouve pas plus que les sections déjà envisagées, associées aux ``places'' de $L_1$ sur $K$.
\margin{on pense ici au cas où $K$ est de type fini sur $\mathbf Q$, bien sûr !}

Mais il en est bien ainsi, comme on voit en considérant $V_1=V\cap L_1$, qui est un anneau de valuation de $L_1$,
\margin{124}
\scanpage{126}
dont le corps résiduel (contenu dans celui de $V$) sera \deleted{alg} fini sur $K$ si celui de $V$ l'est (donc $V_1\ne L_1$) --- donc $V_1$ \deleted{$\longrightarrow$} \inserted{correspond à} une ``place'' du corps de fonctions d'une variable $L_1$ sur $K$. Or, par des fonctorialités essentiellement évidentes, on voit que l'hom $\mathcal E_L^{\overline L}\longrightarrow\mathcal E_{L_1}^{\overline{L_1}}$ envoie $N_V$ dans $N_{V_1}$, $I_V$ dans $I_{V_1}$ --- donc l'image de $N_V$ normalise $I_{V_1}$, donc (si on a une \inserted{germe de} section de $N_V$ sur $\mathcal E_K$) le germe de section obtenu de $\mathcal E_{L_1}$ sur $\mathcal E_K$ normalise $I_{V_1}$ --- et sa restriction à $\mathcal E_K^0$ centralise $I_{V_1}$ ---
\margin{\unsure{Il|?} faut vérifier que le corps des fractions de $V_1$ est bien $L_1$ --- c'est un fait général\ldots}

\uline{Conjecture} Soient $K,L$ des ext. de type fini de $\mathbf Q$, $K\subset L$. Alors

a) Toute section de $\mathcal E_L$ sur $\mathcal E_K$ (ou germe de tel, ça revient au $=$\ldots) normalise un $I_V$ associé à un anneau de valuation \inserted{$V$} de $L$ contenant $K$, à corps résiduel algébrique \inserted{sur $K$} \deleted{\unsure{,}} et \unsure{ce} $V$ est uniquement
\margin{125}
\scanpage{127}
\mspage{63}
\margin{À nouveau se pose la question si $N_V=\operatorname{Norm}(I_V)$. On voudrait que les germes de sections de $\mathcal E_L^{\bar L}$ sur \deleted{$\mathcal E_K$} \inserted{$\mathcal E_{K_1}^{\bar K}$} soient associés \unsure{aux} pts du normalisé $\unsure{\tilde{\mathfrak X}|\tilde{\mathcal Z}}_{L/K}$ de $\unsure{\mathfrak X|\mathcal Z}_{L/K}$ dans $\bar L$ --- qu'elles forment un fibré sur $\unsure{\tilde{\mathfrak X}|\tilde{\mathcal Z}}_{L/K}$, à fibres les $H^1(\deleted{\mathcal E_K}\inserted{\mathcal E_{K_1}}^{\unsure{G|\bar K}},I_V)\ldots$ Ceci est associé à la formule (pour une telle \uline{section} \unsure{$s$} de $\mathcal E_L$ sur $\mathcal E_K$) $\mathcal E_L^{(\mathcal E_K^0)}\simeq I_V$ --- lié à la conjecture que pour $\forall$ \unsure{s-groupe} ouvert de $N_V^0$, ou \unsure{simplt} $\forall$ \unsure{s-groupe} de $N_V$ dont l'image dans $N_V/I_V$ soit \unsure{ouvert}, on a \deleted{\unsure{Centr|tout}} $(\mathcal E_L^{\bar L})^{\unsure{\bar\Gamma|\Gamma}}=I_V$ (ceci pour $\forall$ anneau de valuation $V$ de $L/K$, \unsure{présumé} de corps rés. alg./$K$). En fait, la question est purement géométrique !\ldots}
\unsure{en|à} cette condition --- en fait la section provient d'une section du groupe de décomposition $N_V$ au-dessus de $\mathcal E_K$.

\noindent b) Soit $U$ un modèle ``élémentaire'' de $L$ sur $K$, anabélien. Alors tte section de $\mathcal E_U^{\bar L}$ sur $\mathcal E_K^{\bar K}$ se relève en une section de $\mathcal E_L^{\bar L}$ sur $\mathcal E_K^{\bar K}$.

Elle peut donc se décrire en termes d'un anneau de valuation $V$ de $L/K$, et qd $V$ a un centre sur $U$, \unsure{on n'a qu'à} se \unsure{concentrer} au cas d'un pt $x\in U$ rationnel sur $K$, et de la section définie par $x$. Dans ce cas (\unsure{prenant|posant} tjrs $\Gamma=\mathcal E_K^{\bar K}$) on conjecture encore, bien sûr
\[
(32)\qquad (\mathcal E_L^{\bar L})^{\Gamma^0}=(1)
\]
(et on peut y remplacer $\Gamma^0$ par $\Gamma^0\unsure{\ne1)|\ne)|\subset}\deleted{\unsure{\Gamma}}$
--- mais en fait ceci (en vertu de $\operatorname{Centr}_\Gamma(\Gamma^0)=(1)$) équivaut à la relation analogue \deleted{$\pi_{L\unsure{,|/}K}^{\bar L}$}
\[
(33)\qquad (\pi_{L\unsure{,|/}K}^{\bar L})^{\Gamma^0}=(1)
\]
qui \unsure{se} déduit d'ailleurs, par dévissage,
\margin{126}
\scanpage{128}
de la conjecture analogue déjà faite pour la dim relative $1$.

Pour mieux appréhender cette deuxième partie \inserted{b)} de la conjecture, il faudrait examiner de plus près ce que deviennent, par composition avec $\mathcal E_L^{\bar L}\mathrel{\text{---}}\mathcal E_U^{\bar L}$, les sections de $\mathcal E_L^{\bar L}$ sur $\mathcal E_K$ associées à un $V$, et pour commencer, étudier le sous-groupe \inserted{$\Delta_V$} de $\pi_{U/K}^{\bar L}$ image de $I_V$ --- qui, on le devine, jouera un rôle analogue à un groupe-bout. Dans quelle mesure des $V$ distinctes définissent-elles des $\unsure{\Lambda_V|\Delta_V}$ distinctes ?

Pour étudier cette question, introduisons un modèle $X$ de $L/K$ contenant $U$, $X$ propre, les comp. irréd. de \inserted{$D=$} $X-U$ des diviseurs réguliers \inserted{$D_i$ $(i\in I)$} et $X-U$ à croisements normaux. Pour $x\in\operatorname{supp}D$, considérons le hensélisé \inserted{(non strict)} $\tilde O_x$ de $\underline{O}_{X,x}$, et l'image inverse $\unsure{\bar U_x|\tilde U_x}$ de $U$ dessus. On a
\margin{127}
\scanpage{129}
\mspage{64}
$\bar L$ au-dessus de $\tilde U_x$, bien sûr, et
\[
 (34)\qquad
 1\mathrel{\text{---}}T^{I(x)}\longrightarrow
 \mathcal E_{\tilde U_x}^{\bar L}\longrightarrow
 \mathcal E_{k(x)}^{\overline{k(x)}}\longrightarrow 1
\]
où \deleted{$I(x)$}
\[
 I(x)=\{i\in I\mid x\in D_i\},\qquad T=T_\infty(\bar L^*)
\]
$\overline{k(x)}$ étant la clôture alg. de $k(x)$ induite par un plongement choisi de $\tilde K_x$ (corps des fonctions de $\tilde O_x$) dans $\bar L$, qui fait donc fonction de clôture sép. de $K_x$.

Ceci dit, pour $V$ anneau de valuation de $L/K$, si son centre $x$ \inserted{sur $X$} \unsure{est} $\in\operatorname{Supp}D$, on aura, pour $S=\operatorname{Spec}V$, d'où $\tilde S$ \inserted{(hensélisé)} du pt \unsure{générique|géométrique} $\tilde\eta$ une factorisation de $\tilde\eta\longrightarrow U$ (composé $\tilde\eta\longrightarrow\eta=\operatorname{Spec}(L)\longrightarrow U$) en $\tilde\eta\longrightarrow\tilde U_x$, car $\tilde\eta\mathrel{\text{---}}U$ est induite par $\tilde S\longrightarrow X$ qui se factorise en $\tilde S\longrightarrow\tilde T\longrightarrow X$ ($\tilde T=\operatorname{Spec}\tilde O_x$). Donc \deleted{l'image $N_V$ de} \inserted{l'hom. induit} $\mathcal E_{\tilde\eta}^{L}\longrightarrow\mathcal E_U^{L}$ est \deleted{contenue dans celle de} \inserted{un composé}
\[
 (35)\qquad
 \diagram{\begin{tikzcd}[column sep=large,row sep=small]
 \mathcal E_{\tilde\eta}^{\bar L}\arrow[r]\arrow[d,equal] &
 \mathcal E_{\tilde U_x}^{\bar L}\arrow[r]\arrow[d,equal,"{\unsure{\mathrm{dfn}|\mathrm{déf.}}}"'] &
 \mathcal E_U^{\bar L}\\
 N_V & \deleted{\unsure{N_x}}\ N_x &
 \end{tikzcd}}
\]
(l'image de $\mathcal E_{\tilde U_x}^{\bar L}\longrightarrow\mathcal E_U^{\bar L}$ n'étant d'ailleurs autre que le groupe de décomposition de $x$ dans $\mathcal E_U^{\bar L}$ relatif à $\bar L$\ldots). L'étude de l'hom composé
\margin{128}
\scanpage{130}
se ramène à celle des hom composants.

À ce sujet, sous réserve de vérifications, je présume

$1^\circ)$ que l'hom $\mathcal E_{\tilde U_x}^{\bar L}\longrightarrow\mathcal E_U^{\bar L}$ est \uline{injectif}

--- ou ce qui revient au \unsure{=|même}, que sa partie géométrique
\[
\diagram{\begin{tikzcd}[column sep=large,row sep=small]
\mathcal E_{\tilde U_x/K}^{\bar L}\arrow[r]\arrow[d,equal] & \pi_{U/K}^{\bar L}\\
I_x\simeq T^{(I(x))} &
\end{tikzcd}}
\]
\margin{les deux lettres $I$ ne sont pas les mêmes}
est \uline{injective}.

\margin{C'est visiblement \uline{faux} \inserted{\uline{« en général »}} pour $\operatorname{card}I(x)\ne1$, comme on voit en \unsure{passant} à un autre modèle $X'$ déduit de $X$ en y faisant éclater un pt de $D$\ldots\ Mais y a-t-il un modèle « minimal » $X\supset U$, et pour un tel modèle \unsure{a-t-on} l'injectivité ???}

$2^\circ)$ que \deleted{l'hom} lorsque le corps résiduel de $V$ est alg $/K$ (i.e. $V$ correspondant à \unsure{un pt} de $\unsure{\mathfrak X_{L/K}|\mathcal Z_{L/K}}$ de dim zéro), alors l'hom
\[
 \mathcal E_{\tilde\eta}^{\bar L}=N_V\longrightarrow\mathcal E_{\tilde U_x}^{\bar L}
\]
\deleted{est injectif} \inserted{a une image} \deleted{surjectif} \inserted{ouverte} --- ce qui revient ici (puisque $k(V)$ \deleted{alg} \inserted{fini} sur $k(x)$ donc $\mathcal E_{k(V)}=N_V/I_V\hookrightarrow\mathcal E_{k(x)}$ \inserted{(sous-groupe ouv.)} $\simeq N_x/I_x$), que la partie géométrique
\[
 I_V\longrightarrow I_x=T^{I(x)}
\]
\deleted{est injective} \deleted{\unsure{surjective}} \inserted{a une image ouverte} \deleted{--- \unsure{je pense que} \illegible{3} \unsure{à un sous-groupe ouvert}.} \inserted{--- sans doute est-elle \unsure{=|même} surjective.}

À noter que \unsure{la} question $2^\circ$ est purement
\margin{129}
\scanpage{131}
\mspage{65}
locale autour de $x$ --- elle doit être essentiellement « triviale », que ma présomption s'avère vraie ou fausse --- par contre $1^\circ)$ est une question de nature globale sur $U$, et sans doute d'\deleted{\unsure{autre}} aucune façon triviale.

\deleted{Avant de \unsure{co|ce}} La validité de ces énoncés \deleted{signifierait} \inserted{impliquerait} donc, pour les sections de $\mathcal E_U^{\bar L}$ sur $\mathcal E_K^{\bar L}$ associées à un anneau de valuation de $L/K$, \inserted{de corps résiduel $K$}, que l'image de $\mathcal E_K^{\bar L}$ doit normaliser un \deleted{$I_x$} sous-groupe $\unsure{\Delta_x|\bar\Delta_x}\simeq T^{I(x)}$ de $\pi_{L/K}^{\bar L}$, qui est non trivial si la valuation n'a pas de centre sur $U$, i.e. si la section n'est pas associée à un pt $K$-rationnel de $U$, ce qui est justement dans le sens des conjectures (qui pouvaient d'abord sembler hâtives !) du § $2$.

Avant de me tourner vers l'étude des questions $1^\circ)$ et $2^\circ)$ (et \inserted{aussi} des questions de normalisation et de centralisation soulevées
\margin{130}
\scanpage{132}

précédemment à propos des $N_V,I_V\ldots$), je vais regarder d'un peu plus près l'allure des sous-groupes $I_x\subset\unsure{\pi|\mathcal E}_{L/K}^{\bar L}$, (ou \deleted{\unsure{associés}} plutôt, des classes de \inserted{$\pi$-}conjugaison de tels sous-groupes) associés aux $x\in X\setminus U$, \inserted{pour $x$ variable,} et me \unsure{convaincre} du caractère intrinsèque de ces groupes, indépendamment du choix du modèle $X$ (mais ceci est évident, si on admet ce qu'on vient d'\unsure{énoncer|annoncer} --- que ce sont justement, à sous-groupes ouverts près \deleted{\unsure{tout}} au plus, les images des $I_V$, pour $V\deleted{\unsure{\to}}\inserted{\in}\mathfrak X_{L/K}$ n'ayant \inserted{pas de} centre dans $\deleted{X\setminus U}U\ldots$)

\deleted{Considérons d'abord les}

Soit $J$ une partie de l'ens. fini $I$, et soit
\[
\tag{36}
D_J^{\unsure{\circ}}=\{x\in X\mid I(x)=J\}
=\overbrace{\bigcap_{j\in J}D_j}^{\Delta_J}-\bigcup_{i\in I-J}\Delta_J\cap D_i
\]
\[
D_{\{i\}}=D_i\setminus\left(\bigcup_{j\ne i}\unsure{D_i\cap D_j}\right).
\]
Si $d=\operatorname{card}J$, --- de sorte que l'on voit que $D_J\subset X$ est \inserted{régulier,} de codim $d$ en chacun de ses points (mais bien sûr, il peut être $\varnothing$). Si on appelle ``simplexe'' une partie $J\subset I$

\margin{il vaut mieux définir $D_J$ comme \uline{sous-schéma} \unsure{comme} $\bigcap_{j\in J}D_j$ (36), et définir son \uline{\unsure{support|sous-ouvert}} comme $\{x\in X\mid I(x)\unsure{=|\supseteq}J\}$}

\margin{131}
\scanpage{133}
\mspage{66}

telle que $D_J\ne\varnothing$, \inserted{(i.e. $\unsure{\bar D_J|D_J}=\bigcap_{j\in J}D_j$ \unsure{soit} $\ne\varnothing$)} on trouve sur $I$ une structure de \deleted{\unsure{solide}} \inserted{maquette} simpliciale --- quelle est sa signification topologique? Les \unsure{composantes|comp.} connexes correspondent aux comp. connexes de $D=X\setminus U$ --- ou encore : celles de $\deleted{\unsure{U_\infty}}\hat X\setminus D$, où $\hat X$ est le \inserted{\unsure{l'ombre}} complété de $X$ le long de $D$ --- et $\hat X\setminus D$ a un caractère invariant p.r. à la ``compactification'' choisie $X$ de $U$. S'il était vrai que les $D_J$ \inserted{non vides} sont $\unsure{\infty}$-\unsure{connexes} (ce qui bien sûr n'est \uline{\unsure{pas}} le cas, sauf pour $D_J$ réduit : un point!) on pourrait penser que la topologie de notre \deleted{\unsure{solide}} \inserted{maquette} simpliciale exprime le type d'homotopie de $U_\infty$.

Considérons maintenant le ``localisé'' \deleted{\unsure{de}} (au sens des \unsure{ombres}, ou en un sens étale convenable) de $X$ le long de $D_J$, soit $\deleted{\unsure{\hat X_J}}X_J$, et soit $U_J$ \inserted{$=X_J\setminus D_J$,} la trace de $\unsure{U}$ sur $X_J$. Considérons $\mathcal E_{U_J}$ (de façon précise, $\mathcal E_U^{\bar L}$, si on localise \deleted{p.r. à} étalement), on aura une suite exacte
\[
\tag{37}
\unsure{T^J|\bar T^J}\unsure{\hookrightarrow|\to}\mathcal E_{U_J}^{\bar L}\longrightarrow\mathcal E_{D_J}^{\deleted{\unsure{\tilde D_J}}}\longrightarrow1
\]
--- de façon précise, choisissant un rev. universel $\tilde D_J$ de $D_J$, on trouvera un rev. universel \inserted{(par image inverse)} $\tilde X_J$ de $X_J$, d'où d'ailleurs $\tilde U_J$ et $\tilde U_J\longrightarrow U$ ; \deleted{\unsure{et on}} on désigne par $\tilde L_J$ le corps des fonctions de $\tilde X_J\unsure{\mathrel{\text{---}}|=|\ldots}\tilde U_J$, c'est une ext. \inserted{\unsure{alg.}} de $L$, et

\margin{il vaut mieux travailler avec les ``\unsure{tubulures} étales''\ldots}

\margin{il faut choisir $\tilde D_J$ ``au dessus de $\bar K$''\ldots}

\margin{plutôt $T^J$}

\margin{\uline{NB} Il faut partir de $\unsure{U_{J\bar K}}$, \unsure{si on veut}, d'un rev. universel $\unsure{\tilde V_0}$ de $\unsure{U_{J\alpha}}$ (qu'on va \unsure{multiplier} au dessus de $\bar K$ bien sûr), d'où $\unsure{\tilde X_{J\alpha}}$, $\unsure{\tilde D_{J\alpha}}$.}

\margin{132}
\scanpage{134}

et on choisit un \inserted{$L$-}plongement de $\tilde L_J$ dans $\bar L$.

C'est par cela que (37) est définie --- c'est le \inserted{début de la suite d'homotopie du} fibré, \unsure{moralement}, d'un produit de disques sur $D_J$ (savoir $U_J$ lui-$\unsure{=}$) --- j'ai eu un peu peur d'y mettre un zéro à g, à cause du $\deleted{\unsure{\mathcal E}}$
\[
\pi_2(\bar D_J)\longrightarrow
\underset{\substack{\Vert\\ T^J}}{\pi_1(\text{fibre})}
\qquad (?)!
\]
Mais considérons l'hom
\[
\mathcal E_{U_J}^{\tilde U_J}\longrightarrow\mathcal E_U^{\bar L}
\]
induit par ``l'injection'' $U_J\longrightarrow U$ ; \deleted{elle}\inserted{\deleted{devrait}} heuristiquement on \deleted{devrait avoir}\inserted{a} un diagramme \deleted{induit par} commutatif
\diagram{
\begin{tikzcd}[row sep=large,column sep=large]
U_J \arrow[r] \arrow[d,"\unsure{\varepsilon|i}"'] & U \arrow[d] \\
X_J \arrow[r] & \operatorname{Spec}K=\xi \\
D_J \arrow[u,"\text{équiv. d'homotopie}"] & {}
\end{tikzcd}
}
d'où un hom. de suites exactes d'homotopie, \deleted{relative}\inserted{de celle relative} à la première \deleted{ligne}\inserted{colonne}, dans celle relative à la seconde
\diagram{
\begin{tikzcd}[row sep=huge,column sep=small]
\ldots \arrow[r,dash] &
\overset{\substack{\pi_2(D_J)\\\Vert}}{\pi_2(X_J)} \arrow[r] \arrow[d] &
\overset{\substack{T^J\\\Vert}}{\pi_1(\unsure{\bar G_m^{\times J}})} \arrow[r] \arrow[d] &
\overset{\substack{\mathcal E_{U_J}\\\Vert}}{\pi_1(U_J)} \arrow[r] \arrow[d] &
\overset{\substack{\mathcal E_{D_J}\\\Vert}}{\pi_1(X_J)} \arrow[r] \arrow[d] & 1 \\
\unsure{\ldots|--} \arrow[r,dash] &
\underset{\substack{\Vert\\0}}{\pi_2(\xi)} \arrow[r] &
\underset{\substack{\Vert\\\pi_{U/K}}}{\pi_1(U_{\bar\xi})} \arrow[r] &
\underset{\substack{\Vert\\\mathcal E_U}}{\pi_1(U)} \arrow[r,phantom,"\unsure{-|\to}" description] &
\underset{\substack{\Vert\\\mathcal E_{\unsure{K}}}}{\pi_1(\xi)} \arrow[r] & 1
\end{tikzcd}
}
qui montre que \inserted{le composé} $\pi_2(D_J)\longrightarrow T^J\longrightarrow\pi_{U/K}$ est nul --- donc si on admet que $T^J\unsure{-|\to}\pi_{U/K}$ est injectif, on conclura que $\pi_2(D_J)\unsure{-|\to}T^J$ est nul!

\margin{133}
\scanpage{135}
\mspage{67}

On va donc admettre que (37) définit bien une structure d'\unsure{extensions|extension}, et on aura un hom de suites exactes

\deleted{\unsure{(38)}}
\[
\tag{38}
\diagram{
\begin{tikzcd}[row sep=large,column sep=large]
1 \arrow[r,dash] & T^J \arrow[r] \arrow[d] & \mathcal E_{U_{J,\alpha}}^{\bar L} \arrow[r] \arrow[d] & \mathcal E_{D_{J,\alpha}}^{\tilde D_{J,\alpha}} \arrow[r] \arrow[d] & 1 \\
1 \arrow[r,dash] & \pi_{U/K}^{\bar L} \arrow[r] & \mathcal E_U^{\bar L} \arrow[r] & \mathcal E_K^{\bar K} \arrow[r] & 1
\end{tikzcd}
}
\]
--- où on a rectifié le tir, en faisant attention que (37) n'a de sens telle quelle que si $D_J$ est connexe, \inserted{i.e. $X_J$ connexe} que sinon il faut regarder les composantes connexes $X_{J,\alpha}$ de $X_J$, induisant les composantes \unsure{conn|conn.} $D_{J,\alpha}$ de $D_J$, et $U_{J,\alpha}$ de $U_J$, et écrire la suite exacte pour chaque composante. Bien sûr, il y a des choses plus intrinsèques : écrire avec les topos \deleted{étales \unsure{galoisiens|géométriques}}\inserted{multigaloisiens}, et leurs groupoïdes fondamentaux....\deleted{\unsure{classes de} $\mathcal E_U^{\bar L}$ \unsure{conjugués|conjuguées}}

On voit ainsi que les hom
\[
\tag{39}
T^J\longrightarrow\unsure{\pi|\mathcal E}_{U/K}^{\bar L}
\]
induits par les différents $x\in D_J$ --- \unsure{chacun} n'étant en fait, quand on reprend la définition précise, que défini modulo opérations \unsure{directes|induites} par autom. intérieurs de $\mathcal E_U^{\bar L}$, induits par

\margin{134}
\scanpage{136}

éléments dont l'image dans $\mathcal E_K^{\bar K}$ est dans l'image de $\mathcal E_{k(x)}^{\overline{k(x)}}\unsure{\mathrel{\text{---}}|\longrightarrow}\mathcal E_K^{\bar K}$ --- on voit que pour $x$ variable dans une composante connexe $D_{J,\alpha}$ de $D_J$, sont tous conjugués entre eux dans $\mathcal E_{U/K}^{\unsure{\bar L|\bar K}}$ --- et on peut faire les choix relatifs aux différents $\alpha$ d'une façon \unsure{habile|arbitraire} telle, que maintenant les \unsure{autom.} extérieurs de $\pi_{U/K}^{\bar L}$ qui expriment l'indétermination dans (39) soient ceux provenant des éléments de l'image de
\[
\mathcal E_{D_{J,\alpha}}^{\tilde D_{J,\alpha}}\longrightarrow\mathcal E_K^{\bar K}.
\]

$J=\{i\}$ est particulièrement \unsure{simple}, \deleted{chaque composante \unsure{conn.}}\inserted{dans ce cas d'ailleurs $D_J=D_i$} connexe, et on trouve des hom
\[
\tag{40}
T\xrightarrow{\unsure{\kappa_i|k_i}}\pi_{U/K}^{\bar L},
\]
déterminés a priori mod autom intérieurs près --- donc en fait $\kappa_i$ est un \uline{hom} \uline{extérieur} bien déterminé, qui est \deleted{\unsure{invariant}} \deleted{stable par l'opération} compatible avec les opérations ext. de $\mathcal E_K^{\bar K}$ sur $T$, \uline{et} sur $\pi_{U/K}^{\bar L}$.

\margin{de façon plus précise, le choix d'un rev. universel de $\unsure{\tilde U_x|U_x}$ en implique un de $K$ --- et on suppose que c'est $\bar K$ (avec \unsure{des} \unsure{rigidifications|rigidification}) : i.e. on choisit $\tilde U_x$ au dessus de $\bar K$ --- et alors cela implique un rev. universel de $k(x)$ au dessus de $\bar K$. Ces choix étant faits, on a
\diagram{
\begin{tikzcd}[ampersand replacement=\&]
\mathcal E_{\unsure{\tilde U_x}}^{\unsure{\tilde{\tilde U}_x}} \arrow[r] \& \mathcal E_K^{\bar K}
\end{tikzcd}
}
\ldots}

\margin{Mais a priori ils sont par voie \unsure{géom.} définis mod autom intérieurs de $\uline{\pi_{U/K}^{\bar L}}$, et les opérations \deleted{\unsure{extérieures}} de $\mathcal E_K^{\bar K}$ sur $\pi_{U/K}^{\bar L}$ \unsure{laissent} invariants \unsure{les hom} \inserted{\unsure{extérieurs}} en tenant compte de l'action extérieure ``cyclotomique'' de $\mathcal E_K^{\bar K}$ sur $T^J$\ldots}

\deleted{Si \unsure{on a} $J'\supset J$ \inserted{$\bar D_{J'}\subset\bar D_J$} (i.e. $\unsure{\bar D_{J'}}\deleted{\unsure{\cap D_J}}\ne\varnothing$), choisissons une composante conn. $\unsure{\bar D_{J',\alpha'}}$ de $D_{J'}$, elle est contenue dans une comp. conn.}

\margin{135}
